\paragraph{Scope and validation.}
\label{sec:design-discussion}
Sparse integration needs native selection indices, recoverable selector
state and all attention branches, physical-block mapping, and the directory/
lease interface above. Poor predictions can use exact demand fetch or
full-layer loading; dense attention uses the full historical set.
Standalone requests use zero slack. Prefill failures require profiles for
the lost chunk, not a universal 64-token tail.

The present evidence consists of native selection traces, attributed
prediction results, and a single-request pipeline model using measured
forward latency and assumed transfer bandwidth. It does
not validate the depth-score profiles, failure-risk calibration, checkpoint
sensitivity, active hot-set migration, concurrent retry fairness, or
foreground interference. End-to-end tests must include recovery during
demotion, stale-location races, DRAM exhaustion, copy backlogs, and shared
bandwidth pressure. The specified interfaces and fallbacks make those
claims testable; they do not establish performance or a completed implementation.
Validation must also replay common source failures through the dependency
DAG, compare the local policy with feasible small-DAG allocations, and
report regret, including zero or negative Peer benefit. Native selector
traces and the measured batched forward do not by themselves validate
recovery equivalence to token-wise decoding.
